\documentclass[11pt]{article}
\usepackage{amsmath,amssymb}
\usepackage[margin=1in]{geometry}
\usepackage{hyperref}

\title{Ideas Parked --- Not Yet Merged}
\author{Staging document; each entry marks its eventual destination}
\date{}

\begin{document}
\maketitle

\noindent This document holds ideas discussed against Pattern Arithmetic
v15\_19, Inference, and Assembly, none of which have been written into
any of those three files yet. Each entry states its status and its
intended destination. Nothing here is settled; the point of this file
is to let that be true without leaving Paper 1 mid-edit.

\section{Destined for Pattern Arithmetic (untouched, remains open ---
now far more precisely characterized)}

\paragraph{Mix across differing directions.} Still open. No formula
has changed, no route has been chosen. What changed is how precisely
the gap itself can now be stated --- an extended discussion worked
through what Mix actually \emph{is}, why several everyday-seeming
counterexamples turn out not to be counterexamples, and split the
open question into two genuinely separate parts that had been sitting
together as one.

\textbf{[The reading, argued but not yet adopted as the paper's own
position.]} Mix is not usefully described as ``multiple vectors
becoming one vector.'' That framing is general enough to make the
gap look like an oversight. A better reading: Mix takes several
reports already agreed to be answers to \emph{one} shared question,
and returns a verdict about that question. The shared $\hat n$ isn't
a boundary condition sitting on top of a more general operator --- it
is what makes ``these are answering the same thing'' true in the
first place. In Mix's already-defined domain, $\hat n$ itself was
never plural; only $\chi$ ever varied.

\textbf{[Supporting evidence for that reading, already on file
elsewhere in this project.]}
\begin{itemize}
\item Same-direction, opposite-phase, equal strength already returns
      the null state, black --- not a fabricated compromise.
      Genuine contradiction is allowed to cancel to nothing, which
      only makes sense if agreement and disagreement were
      well-defined to begin with, which requires one shared axis.
\item Unmixing is provably, structurally ill-posed (verified: two
      genuinely different input pairs land on an identical output).
      Metamerism was checked as independent real-world precedent for
      the same fact: many-to-one perceptual reduction is genuinely
      lossy, not a design flaw to route around.
\item The paper's own opening line already states the principle
      directly: ``Two notes yield a chord, not a larger integer... the
      word `add' names a different rule.'' A chord is several notes
      sounding together \emph{because} they already share one frame
      --- a key, a scale. Not a generic combination of any two
      sounds.
\end{itemize}

\textbf{[Resolved, not open: the ``we mix all the time'' challenge.]}
Four everyday examples were checked directly against this reading.
Sugar dissolving in tea is one dial (sweetness) turning further along
itself --- never needed differing directions. Red plus green
additively is already confirmed, in the paper's own words, to sit
``at the same place'' --- same $\hat n$, different $\chi$ --- already
inside Mix's settled domain; ``a new reality'' (yellow) is a new
position within one shared phase-space, not evidence that different
directions were combined. Layer-based paint software was checked
directly and confirms the framework's own distinction under different
names: compositing separate layers is non-destructive and reversible
(a bowl, $\oplus$, nothing discarded), while flattening is
documented, across every source checked, as explicitly destructive
and irreversible (Mix, $\boxplus$) --- ``once flattened, individual
layers cannot be separated again.'' ``Unmerge them later'' is only
ever possible because the layers were never actually flattened; the
moment they genuinely are, every source agrees: gone, matching
unmixing's own ill-posedness exactly. The one example that is
\emph{not} resolved --- representing a whole painting (color,
brushstroke, texture) as a single vector --- turns out to be Gap 1
itself, restated as a question rather than evidence that it is
already answered.

\textbf{[New: the open question splits into two, not one.]} Scope
(different sorts entirely, like color versus sweetness) was already
handled --- route 1, refuse outright, a category error rather than a
hard case. What remains splits further, for the case of one shared
sort whose members sit at genuinely different positions on one
sphere (rain and drought sharing a weather-concept sphere, the way
red and green share a color one):
\begin{enumerate}
\item \emph{Embedding validity.} Does geometric closeness on that
      shared sphere actually track real semantic agreement or
      disagreement at all? For hue, yes, by deliberate construction:
      the circle is built so that maximally different hues sit
      antipodally. For an arbitrary sphere of concepts, nobody has
      stated this. It is not given for free by placing two things on
      one sphere.
\item \emph{The merge rule.} Only given a validated embedding does
      the original question --- what should combining two positions
      produce --- become well-posed at all. This is the question the
      existing paragraph already names; it was previously the whole
      of Gap 1, and is now understood to be its second half.
\end{enumerate}

\textbf{[Still confirmed not to help, checked directly, carried
forward from prior sessions.]} Magnitude behaving well (self-bounding
near genuine opposition, verified via the law of cosines) does not
license the direction that comes with it --- a small, well-behaved
magnitude can still carry an illegitimately-produced direction. The
Hill form and Assembly's Hill bijection both fix magnitude
saturation only, verified twice, never direction. Bowls, per-channel
scoring, space-time gating, and the candlestick merge rule are all
genuine, useful ways of \emph{not needing} a combined direction ---
none of them state what one would be.

\textbf{[Open, unchanged in substance, changed in precision.]} No
merge rule exists for either sub-question above. The three routes on
file --- refuse outright, rotate onto a shared axis via a stated $T$
first, or build a genuine two-stage operator --- remain unchosen. The
result of any eventual resolution is deterministic and reproducible,
not random; what is missing is licensing, not predictability.

\section{Destined for Pattern Arithmetic (advances an existing open item)}

\paragraph{Location as a sort --- a third candidate.} The existing
paragraph names two candidates and adopts neither. A third was tested
here: dedicate separate, disjoint dimensions to content, to time, and
to location, rather than letting any of them share an axis.

\textbf{[Verified.]} Two shifts acting through disjoint content
dimensions commute exactly (checked numerically, not approximately:
time-then-location and location-then-time landed on the identical
point to full floating-point precision). Two shifts forced to share a
single content dimension do not commute, confirming the disjointness
is what matters, not rotation itself.

\textbf{[Verified.]} Content stays recoverable after both shifts:
projecting back onto the content dimensions alone reproduced the
original direction closely, while the time-axis and location-axis
each carried a clean, separately-readable amount of shift, with
nothing blended between them.

\textbf{[Open, unchanged.]} The encoding itself --- what specific
rotation a given time value or location corresponds to --- is not
stated anywhere. The geometry now supports it cleanly; the mapping
does not exist yet.

This is real, checkable content, not yet written as an addition to
the existing paragraph.

\section{Destined for Pattern Arithmetic (open item stays open, unchanged)}

\paragraph{Reducing a path by merging close actors.} The existing
paragraph asks for three things: a stated distance, a merge rule, and
a bound on how far a reduced path's endpoint may drift from the
unreduced one. Discussion here supplied a distance (see below) and
partially addressed the merge rule, but only for magnitude, not for
direction --- and direction is the harder half this item is actually
about.

\textbf{[Verified.]} Switching the strength calculation from the
companion paper's clip to Inference's already-established Hill form,
$r'=|z|/(K+|z|)$, changes nothing about the resulting direction
$\chi'=\arg(z)$. Checked directly on two constructed sequences with
identical content and opposite order: both gave the same angle under
either formula. The Hill substitution fixes evidence-count blindness
in strength; it does not touch the order-blindness this open item is
actually asking about.

\textbf{[Open, unchanged.]} No merge rule for direction exists. No
drift bound exists for either magnitude or direction. This item
remains exactly as open as it was before this discussion, now with
the magnitude half more precisely separated from the direction half.

\section{Destined for Inference (new mechanism)}

\paragraph{Candlestick summary: a richer merge rule for reference-atom
building.} Not a resolution of the Paper 1 item above --- a different
question. Paper 1's open item is about \emph{walking}: can a shortened
path still be walked to roughly the same endpoint. This is about
\emph{reading}: how to summarize a group of atoms for comparison
against a reference atom, which is Inference's own established
business (its reference-atom section), not $T$'s.

\textbf{[Verified.]} Open/close/high/low, computed the way a trading
candle is, distinguishes exactly the case flat Mix could not. Two
sequences, identical content, opposite order --- naive Mix gave the
identical angle ($47.5°$) for both. The candle's open and close swapped
between them ($20°\!\to\!160°$ versus $160°\!\to\!20°$), correctly
reading one as ending badly and the other as ending well. High and low
stayed identical for both, correctly capturing that only the
\emph{range} was shared, not the \emph{direction of travel}.

\textbf{[Settled, matches existing machinery.]} The scope rule ---
every atom belongs to exactly one time unit, no overlap between
buckets --- is not new; it is precisely the discrete-time-bucket
metadata already defined on a walk's own steps in the companion paper.

\textbf{[Open.]} This gives Inference a genuinely better reference-atom
merge rule than flat Mix. It is not proposed as a fix for Paper 1's
walking-specific open item, which remains untouched above. Whether
open/close/high/low itself needs a coherence check analogous to
$\bar R_c$ has not been considered.

\section{Destined for Inference (naming caution, not yet a mechanism)}

\paragraph{``Inference Mandelbrot'' is not Assembly's Mandelbrot, and
should not share its name.} Assembly's multi-zoom promotes an
unordered thing (a settled camp) and is fully handled by the Hill
bijection already verified there. Whatever would promote an
\emph{ordered} thing (a walked path, a movie, a time series) into one
higher-level atom is a different, harder problem --- it is the
direction half of the path-reduction item above, restated from a
different angle. \textbf{[Naming flag only.]} If this is ever built,
it wants its own name; reusing ``Mandelbrot'' for it would be the same
kind of collision already caught and corrected for ``collapse,''
``singularity,'' and the quantum swap test.

\section{Destined across three papers (new, split by piece)}

\paragraph{Meta vectors: a second kind of vector alongside information
vectors.} Proposed: an information atom may carry two distinct kinds
of vector. Information vectors are Axiom 7's bundle $v_i$'s, capped at
$r_i\le1$, potentially several sorts $\tau$ across a whole record.
Meta vectors describe the atom rather than being part of its content
--- three named so far: $V_{\text{derived}}$, link vectors (pointers
between atoms, one- or two-directional), and a timestamp vector.
Explicitly left open whether more meta-vector types are needed. Meta
vectors are not capped at 1, and grouping several atoms' meta vectors
together on like terms is proposed as a distinct mechanism for
building higher structures. This is one idea that lands in three
different places once traced to where each piece actually belongs.

\textbf{[Destined for Pattern Arithmetic.]} The taxonomy itself ---
information vectors versus meta vectors as a formal, named distinction
sitting on top of Axiom 7 --- is structural, the same kind of addition
Axiom 7 itself was. $V_{\text{derived}}$ already lives there. A
timestamp vector, if adopted, would be a third, genuinely different
treatment of time from the two already on file above: distinct from
the discrete-bucket metadata already defined on a walk's own step, and
distinct from the dedicated-axis encoding in the location-as-a-sort
entry. Not yet clear whether more than one of the three should
coexist, or which.

\textbf{[Destined for Assembly.]} Link vectors specifically. This is
the mechanism the shelf's own neighbour-graph note was already naming
without a concrete form: ``a data structure on seats after rest, not
$T$.'' Grouping link vectors on like terms is collecting edges into a
graph, not merging them --- gathering, in the $\oplus$ sense --- so
it carries none of the open questions below it.

\textbf{[Destined for Inference.]} The grouping mechanism itself:
querying several atoms' meta vectors, confirming they sit close
together, and promoting one shared meta-vector up to the group as a
whole. This is the space-time gating idea from earlier discussion,
generalized into a standing principle. Timestamp and location get
\emph{used} here; they are only \emph{defined} in Pattern Arithmetic
above.

\textbf{[Verified, carried forward as a caveat rather than a new
finding.]} Combining several atoms' $V_{\text{derived}}$'s stays safe
only under the constraint the original already carries: auxiliary,
never promoted to a trusted, walkable atom on its own. Checked
directly that neither Inference's Hill form nor Assembly's Hill
bijection touches this constraint --- both fix magnitude only, never
direction --- so grouping $V_{\text{derived}}$'s inherits that same
limit rather than escaping it.

\textbf{[Note added here.]} Link vectors were separately discussed as
one possible mechanism for a limited, honest form of ``uncollapse'':
not inverting Mix (still impossible, unaffected by anything here),
but keeping an explicit, external pointer from a collapsed atom back
to the bowl that produced it, created at collapse time, before the
bowl itself might otherwise be discarded. This is navigation via a
kept record, not reconstruction from the collapsed result --- the
distinction the white-light/prism discussion turned on, confirmed
there to matter.

\section{Considered and found to need nothing filed}

For completeness: the reference-atom and precursor-witness mechanism,
the bundle-vector / uncapped $V_{\text{derived}}$ distinction from
Mix, and the $\pi=(x;w_1,\ldots,w_n)$ poem trajectory were all
discussed at length and all turned out to already be fully present in
Inference and Paper 1 respectively. Nothing from those threads is
parked here, since there was nothing left unresolved to park.

\end{document}
